% !TEX root = AImath.v4.tex
\chapter*{ 前言}
%\markboth{前言}{前言}

\addcontentsline{toc}{chapter}{前言}

\begin{introduction}

\item 现在看起来的“显而易见”，在当年却是“困难重重”。这句话深刻地概括了人工智能发展的历史轨迹，也揭示了人类智慧演进的根本规律。今天我们习以为常的技术成就，在其诞生之初都曾面临着巨大的挑战和质疑。

\end{introduction}

\section{从“困难重重”到“显而易见”：智能发展的历史辩证法}

现在看起来的“显而易见”，在当年却是“困难重重”。

这句话，也许是理解人工智能发展历程的最好钥匙。今天我们轻松使用语音助手、个性化推荐、智能搜索摘要、实时翻译，甚至让大模型生成文本、代码与图像，驱动办公 Copilot 与多模态对话——很难想象这些能力背后经历过怎样的艰难探索。每一个今天看来理所当然的功能，都是无数研究者在黑暗中摸索、在失败中坚持、在质疑中前行的结果。

人类历史上，每一次重大的技术突破都遵循着相似的轨迹：从被认为不可能，到被视为困难，再到变得显而易见，最终成为日常生活的一部分。这种认知转变的背后，隐藏着科学进步的深层逻辑。

在古代，人们认为飞行是神的特权，人类永远无法像鸟儿一样翱翔天空。莱特兄弟的第一次飞行被许多人视为杂技表演，而非真正的技术突破。然而，仅仅几十年后，航空技术就彻底改变了人类的生活方式。今天，我们把乘坐飞机视为最普通不过的出行方式。

人工智能的发展历程同样如此。1950年，当艾伦·图灵提出“机器能思考吗？”这个问题时，大多数人认为这是科幻小说的幻想。机器怎么可能具有人类独有的思维能力？然而，从1997年深蓝击败国际象棋世界冠军卡斯帕罗夫，到2016年AlphaGo征服围棋巅峰，再到2022年ChatGPT展现出令人惊叹的语言理解能力，我们见证了一个又一个“不可能”变成“显而易见”。

在古代的神话传说中，技艺高超的工匠就可以制作人造“人”，并为其赋予智能或意识。希腊神话中的塔罗斯（Talos）是由火神赫菲斯托斯打造的青铜巨人，负责守卫克里特岛；犹太传说中的泥人哥连（Golem）由拉比用泥土塑造，通过神秘的仪式获得生命。这些古老的想象，预示着人类对创造智能的永恒渴望。

要理解这种“从困难到显而易见”的逻辑，最清晰的路径，是回到数学与计算的演化主线。

\section{数学史：智能追求的发展主线}

本质上，人类发展史就是一部智能简史。从人类文明最早出现数的概念，到数的标记和运算，从算术到代数、几何、抽象代数、拓扑、群论等越来越复杂和专门的数学分支，数学史的发展主线背后，是对智能的极致追求。

从最初的计数开始，人类就在不断追求计算越来越高效、越来越自动化、越来越智能。算盘的发明让计算摆脱了手指的限制；对数表的出现使复杂运算变得简单；机械计算器的诞生标志着自动化计算的开端。每一步进展，都是在拓展人类智能的边界。从人脑和手算发展到一定程度，对更高效计算效率的追求必然导致计算机的发明，进而催生人工智能。这不是偶然，而是历史发展的必然逻辑。


“语言是思想的边界，技术是思想的实现。”这句话深刻地揭示了人工智能发展的哲学内核。人类的每一次技术突破，都是在拓展思想的边界，而技术的进步又反过来推动思想的深化。

语言不仅是交流的工具，更是思维的载体。当我们发明新的概念、创造新的词汇时，实际上是在扩展思维的可能性。计算机科学的发展历程完美地诠释了这一点：从“算法”到“程序”，从“数据结构”到“人工智能”，每一个新概念的诞生都标志着思维边界的拓展。

技术则是思想的物化。当抽象的数学理论转化为具体的算法，当理论模型变成可运行的程序，思想就获得了改变世界的力量。深度学习的发展历程就是这种转化的典型例子：从1943年麦卡洛克和皮茨提出人工神经元的数学模型，到1986年反向传播算法的完善，再到2012年深度卷积网络在图像识别上的突破，理论与实践的相互促进推动了整个领域的飞跃。

数学史的发展主线背后，隐藏着人类对智能的极致追求。从最基础的算术运算，到复杂的抽象代数理论，每一个数学分支的诞生都是为了解决更复杂的问题，实现更高层次的智能。

正如歌德把历史称作“上帝的神秘作坊”，在这个“作坊”里，堆积着对人类来说看似无关紧要的小事。只在极少时候，它们才会被某些时刻点亮，史蒂芬·茨威格称之为“人类的星光时刻”。数学发展史上的每一个重大突破，都是这样的“宿命时刻”——一连串的事件被紧密地编织在一起，经过漫长的发展，然后突然成熟。

欧几里得几何的公理化体系，牛顿微积分的发明，非欧几何的发现，集合论的建立，哥德尔不完备定理的证明——这些“宿命时刻”如繁星般“熠熠生辉、永不熄灭”。它们在数百年的黑暗中闪耀，用数学照亮了人类的历程。

这条主线直接推动了智能层次的跃迁——从模仿到创造。

\section{智能的层次：从模仿到创造}

恰以其分地拿捏尺度，显然是智能的外在表现之一。“过犹不及”说得就是这个理。智能不仅仅是计算能力的体现，更是对复杂情况的精准判断和适度响应。

在人工智能的发展过程中，我们可以观察到智能的不同层次。最初的人工智能系统主要依靠规则和逻辑推理，这种“符号主义”的方法虽然在某些领域取得了成功，但面对复杂的现实世界问题时显得力不从心。

随着机器学习的兴起，人工智能开始具备从数据中学习的能力。这种“连接主义”的方法让机器能够识别模式、进行预测，但本质上仍然是对已有数据的统计分析和模仿。

深度学习的突破标志着人工智能向更高层次的跃升。通过多层神经网络的复杂变换，机器开始展现出某种“理解”能力，能够处理自然语言、识别图像、甚至进行创作。然而，这种智能仍然主要基于模式匹配和统计关联。

真正的挑战在于如何让机器具备创造性思维和抽象推理能力。这需要我们重新思考智能的本质，探索认知科学、神经科学、哲学等多个领域的交叉融合。

\section{数学：智能的基石与必然逻辑}

数学，是人工智能的底层结构与连接逻辑。正如史蒂夫·乔布斯所言，connecting the dots 只能在回望时发生——当我们把“点”连起来，才看见线性代数、概率论、微积分、信息论等学科如何共同托举起计算与智能的进化。

每一个看似突然的技术突破，背后都有着深厚的历史积淀。计算机的发明离不开数学、物理学、工程学等多个学科的发展；人工智能的兴起建立在计算机科学、认知科学、神经科学等领域的基础之上。

这种必然性并不意味着发展过程是平坦的。相反，每一次重大突破都伴随着激烈的争论、反复的挫折、长期的沉寂。深度学习就经历了三次起伏：1940-1960年代的感知机热潮，1980-1990年代的反向传播复兴，以及2010年代至今的深度学习革命。每一次低谷都被认为是“AI冬天”，每一次复兴都被视为“奇迹”。

数学在人工智能发展中扮演着至关重要的角色。从最初的逻辑推理到现代的深度学习，每一次重大进展都离不开数学理论的支撑。线性代数为神经网络提供了计算框架，概率论为机器学习奠定了理论基础，微积分使得梯度下降成为可能，信息论指导着数据压缩和特征提取。

这些数学工具的发展历程同样充满了“困难重重”到“显而易见”的转变。牛顿和莱布尼兹发明微积分时，面临的不仅是技术难题，更有来自学术界的激烈争议。今天，微积分已成为理工科学生的必修课程，其重要性不言而喻。

概率论的发展更加曲折。从帕斯卡和费马研究赌博问题开始，到贝叶斯定理的提出，再到现代统计学的建立，概率论经历了从“数学游戏”到“科学基础”的转变。今天，贝叶斯方法已成为机器学习的核心工具之一。

\section{本书的目标与挑战}

本书试图通过数学的视角，揭示人工智能发展的内在逻辑。我们将探讨那些看似“显而易见”的技术背后的数学原理，回顾那些曾经“困难重重”的问题是如何被一步步解决的。

这是一个充满挑战的任务。人工智能是一个高度跨学科的领域，涉及数学、计算机科学、认知科学、神经科学、哲学等多个学科。要在有限的篇幅内既保持数学的严谨性，又确保内容的可读性，需要在深度和广度之间找到平衡。

我们的目标不仅是传授知识，更是启发思考。希望读者能够理解，今天的“显而易见”正是昨天无数研究者“困难重重”探索的结果，而明天的突破，或许就隐藏在今天看似不可能的想法之中。

正如伏尔泰在《老实人》中所说：“如果没有这场灾难，事物就无法继续进行了。因为一切都是最好的安排。”每一次挫折都是为了更大的突破做准备，每一个“困难重重”的今天都孕育着“显而易见”的明天。

在这个人工智能快速发展的时代，我们既要保持对技术进步的敬畏，也要对未来的挑战保持清醒的认识。数学作为人工智能的基石，将继续指引我们探索智能的奥秘，推动人类文明向更高层次发展。

让我们带着这样的认识，开始这段探索人工智能数学基础的旅程。在这个旅程中，我们将见证“困难重重”如何转化为“显而易见”，体验数学之美如何照亮智能之路。

\section{致读者}

这本书不是一本传统的教科书，也不是一本纯粹的学术专著。它更像是一次思想的漫游，一场关于智能本质的哲学思辨，一段追溯数学与人工智能关系的历史探索。

我们将从最基础的数学概念出发，逐步深入到人工智能的核心算法。每一个公式背后，都有着深刻的思想；每一个算法背后，都有着动人的故事。我们不仅要理解“是什么”和“怎么做”，还要思考“为什么”和“意味着什么”。

在这个人工智能技术日新月异的时代，保持对基础理论的敬畏和对未来发展的思考，或许比掌握具体的技术细节更为重要。因为技术会过时，但思想永恒；工具会更新，但智慧长存。

愿这本书能够成为你探索智能奥秘的指南针，在“困难重重”的学习路上为你点亮一盏明灯，最终让那些看似高深的理论变得“显而易见”。


\newpage